\Needspace{6\baselineskip}
\begin{proof}[Proof of Proposition~\ref{prop:rewrite-competitive-equilibrium}]
\phantomsection
\label{proof:rewrite-competitive-equilibrium}
Free entry and the common unit cost $1/B$ imply $p_X=1/B$.
Because improvements are immediately available to every supplier and
research receives no separate payment, a supplier's net profit is $-M$.
Its optimal research expenditure is zero. The RSI technology
\eqref{eq:rewrite-capability-law} then gives $B=B_0$, establishing
equation~\eqref{eq:rewrite-competitive-no-research}. The remaining
argument constructs and verifies the allocation at this fixed efficiency.

\emph{Net production and normalized dynamics.}
Write $k=K/(AL)$, $c=C/(AL)$, and $f(k)=F_{B_0}(k,1)$.
For $k>0$, the objective in
\eqref{eq:rewrite-competitive-net-production} is strictly concave in $U$.
Its marginal product tends to infinity at zero and to zero at infinity,
so there is a unique positive maximizer. The first-order condition is
$U=(1-\alpha)s_XY$. In normalized units, write
$u=U/(AL)$ and $y=Y/(AL)$. Differentiating this condition, using
\eqref{eq:rewrite-inverse-demand-elasticity}, gives
\[
 \frac{\dd\log u}{\dd\log k}=\frac{\alpha}{e_X},\qquad
 f'(k)=\alpha\frac{y}{k},\qquad
 \frac{kf''(k)}{f'(k)}
 =-\frac{(1-\alpha)(1-s_X)}{\sigma e_X}<0.
\]
Thus $f$ is smooth, increasing, and strictly concave, with $f(0)=0$
and $f'(0+)=\infty$. As $k\to\infty$, the optimal $u$ tends to
infinity and $s_X\to1$. The pricing condition gives
$X/Y\to B_0(1-\alpha)$, while
\eqref{eq:rewrite-composite-ai-limit} gives
$Z/X\to\omega_X^{\sigma/(\sigma-1)}$. It follows that
\begin{equation}
 \lim_{k\to\infty}f'(k)=\lim_{k\to\infty}\frac{f(k)}k
 =\mathcal R^{\mathrm{mc}}(B_0)+\delta,
 \qquad
 f(k)-kf'(k)=(1-\alpha)(1-s_X)y>0.
 \label{eq:rewrite-competitive-net-limits}
\end{equation}
This also derives equation~\eqref{eq:rewrite-marginal-cost-return}.

Set $d=\rho-n>0$ and $G(k)=f(k)-(\delta+n+\gamma)k$.
The household's objective, up to a positive factor and an additive
constant, is $\int_0^\infty e^{-dt}\log c_t\dd t$.
Equations~\eqref{eq:rewrite-competitive-dynamics} become
\begin{equation}
 \dot k=G(k)-c,\qquad
 \frac{\dot c}{c}=G'(k)-d,\qquad
 \lim_{t\to\infty}e^{-dt}\frac{k_t}{c_t}=0.
 \label{eq:rewrite-competitive-normalized}
\end{equation}
An identity useful for constructing the trajectory is
\begin{equation}
 \frac{\dd}{\dd t}\log\frac ck
 =\frac ck-d-\frac{f(k)-kf'(k)}k.
 \label{eq:rewrite-competitive-consumption-ratio}
\end{equation}

\emph{Efficiency below the threshold.}
If $\mathcal R^{\mathrm{mc}}(B_0)<\rho+\gamma$, there is a unique
$\bar k>0$ satisfying $G'(\bar k)=d$. The stationary consumption
level is $\bar c=G(\bar k)>d\bar k>0$, by
\eqref{eq:rewrite-competitive-net-limits}. The Jacobian at this point is
\[
 \begin{pmatrix}d&-1\\ \bar c f''(\bar k)&0\end{pmatrix},
\]
whose determinant is negative. Its stable curve has positive slope,
with $\dot k,\dot c>0$ on its left branch and
$\dot k,\dot c<0$ on its right branch.

Both branches extend to every $k>0$. On the left branch,
$dk<c<G(k)$; on the right, $c>\max\{dk,G(k)\}$.
To verify these restrictions, at $c=dk$ equation
\eqref{eq:rewrite-competitive-consumption-ratio} points toward
$c/k<d$, from which the path cannot return to $\bar c/\bar k>d$.
At $c=G(k)$, the derivative of $\dot k=G(k)-c$ is
$-\dot c$. Below $\bar k$, crossing this curve makes capital decrease
while consumption increases, precluding convergence to $\bar k$.
Above $\bar k$, it makes capital increase while consumption decreases,
again precluding convergence. The stable branches cannot make either
crossing.

Continuing the left branch backward therefore takes $k$ to zero:
otherwise its bounded, monotone coordinates would approach a positive
stationary point different from $(\bar k,\bar c)$, which does not exist.
The right branch extends to arbitrarily large $k$. It cannot approach
another finite stationary point or diverge in consumption at a finite
positive $k$: along this branch
$\dd c/\dd k=c(G'(k)-d)/(G(k)-c)$ remains bounded when $c$ becomes
large on a compact interval of $k$.
For each $k_0>0$, the stable curve consequently supplies a positive
trajectory converging to $(\bar k,\bar c)$ and satisfying the
transversality condition in \eqref{eq:rewrite-competitive-normalized}.

\emph{Efficiency at or above the threshold.}
If $\mathcal R^{\mathrm{mc}}(B_0)\geq\rho+\gamma$, strict concavity
implies $G'(k)>d$ and $G(k)>dk$ at every finite $k>0$.
For any $k_0>0$, consider initial consumption in
$[dk_0,G(k_0)]$ and the region
\[
 k>0,\qquad dk<c<G(k).
\]
At its lower consumption boundary,
\eqref{eq:rewrite-competitive-consumption-ratio} points strictly out
of the region. At its upper boundary, $\dot k=0$ and $\dot c>0$,
so the vector field also points strictly out. Inside, both $k$ and $c$
increase. There is no finite-time escape to infinity because concavity
bounds $G(k)$ above by an affine function. Continuous dependence on
initial consumption makes the sets of paths that leave through the
two boundaries open and disjoint. Both sets are nonempty, as shown
by the endpoint initial choices. They cannot exhaust the connected
interval of initial choices. Hence at least one trajectory remains
in the region forever.

Along this trajectory $k\to\infty$. A finite limit would also bound
$c$ and require a positive stationary point, but $G'(k)>d$ excludes one.
Moreover, $d<c/k<G(k)/k$, and the upper bound satisfies
\[
 \frac{G(k)}k\longrightarrow\mathcal R^{\mathrm{mc}}(B_0)-n-\gamma.
\]
If $\mathcal R^{\mathrm{mc}}(B_0)=\rho+\gamma$, these bounds imply
$c/k\to d$. If the inequality is strict, the same limit follows
from \eqref{eq:rewrite-competitive-consumption-ratio}:
$(f-kf')/k\to0$, and any persistent excess of $c/k$ over $d$
would make the ratio increase until it violated its bounded upper
limit. More precisely, for any $\varepsilon>0$ and sufficiently large
$t$, reaching $c/k\geq d+\varepsilon$ implies
$\dd\log(c/k)/\dd t\geq\varepsilon/2$ thereafter until that
contradiction occurs. Thus
\begin{equation}
 \frac ck\to d,\qquad
 \frac{\dot k}{k}\to\mathcal R^{\mathrm{mc}}(B_0)-\rho-\gamma.
 \label{eq:rewrite-competitive-capital-growth}
\end{equation}
The transversality condition holds because $c/k>d$.
At equality, capital per effective worker diverges but its growth rate
tends to zero.

\emph{Optimality, uniqueness, and decentralization.}
The preceding constructions satisfy global optimality, not only the
Euler equation. Let $\lambda_t=e^{-dt}/c_t$ along a constructed path;
then $\dot\lambda=-G'(k)\lambda$.
For any alternative feasible path $(\widetilde k,\widetilde c)$
with the same $k_0$, concavity gives
\[
 \int_0^T e^{-dt}(\log\widetilde c-\log c)\dd t
 \leq-\lambda_T(\widetilde k_T-k_T)
 \leq\lambda_T k_T\longrightarrow0.
\]
The constructed path has finite utility: in the first case $c$ converges
to a positive constant; in the other cases it is positive, increasing,
and at most exponentially growing. Thus the inequality verifies
optimality in the fixed-$B_0$ allocation problem. Strict concavity of
log utility and concavity of the resource constraint imply uniqueness
of consumption; the state equation then gives uniqueness of capital.
This verification does not allow a planner to improve $B$: it concerns
the allocation with the technology fixed by the private research choices.

The maximizing $U$ in \eqref{eq:rewrite-competitive-net-production},
the factor prices in \eqref{eq:rewrite-capital-return} and
\eqref{eq:rewrite-wage}, and $p_X=1/B_0$ decentralize this allocation.
In particular, $wN=(AL)[f(k)-kf'(k)]$, so the household budget coincides
with the resource constraint. Its Euler and transversality conditions
verify its optimum at these prices. Strict concavity also implies that
any other competitive equilibrium would solve the same fixed-technology
allocation problem and hence have the same aggregate allocation.

\emph{Growth rates and income shares.}
Below the threshold, convergence to $(\bar k,\bar c)$ gives
$r\to\rho+\gamma$, $g_{Y/N},g_w\to\gamma$, and a positive limiting
labor share, because optimal $u$ and $s_X$ have finite interior limits.
At or above the threshold, $k,u\to\infty$ and $s_X\to1$.
Differentiating the static allocation gives
\[
 \frac{\dd\log y}{\dd\log k}
 =1-\frac{(1-\alpha)(1-s_X)}{\sigma e_X}\longrightarrow1,
 \qquad
 \frac{\dd\log(w/A)}{\dd\log k}
 =\frac{\alpha}{\sigma e_X}\longrightarrow\frac1\sigma.
\]
Combining these elasticities with
\eqref{eq:rewrite-competitive-capital-growth} proves the stated growth
limits, including the equality case. Equation
\eqref{eq:rewrite-competitive-net-limits} gives the interest-rate limit,
and $wL/Y=(1-\alpha)(1-s_X)\to0$ completes the proof.
\end{proof}
